% =====================================================================
%  Model Scanner (Jorak) — Détection reference-free de l'abliteration
%  VERSION SIMPLIFIÉE (français). Justification mathématique condensée.
%
%  COMPILATION : Overleaf -> moteur pdfLaTeX.
%  FIGURES     : les 3 PNG sont dans docs/figures/ :
%                  scatter.png, heatmap-ablated.png, heatmap-censored.png
%
%  Version anglaise : docs/report_en.tex (maintenue en parallèle avec ce fichier)
% =====================================================================
\documentclass[11pt,a4paper]{article}

\usepackage[T1]{fontenc}
\usepackage[utf8]{inputenc}
\usepackage[french]{babel}
\usepackage{lmodern}
\usepackage{amsmath,amssymb,amsthm,mathtools}
\usepackage{geometry}
\geometry{margin=2.4cm}
\usepackage{booktabs}
\usepackage{array}
\usepackage{graphicx}
\usepackage{float}
\usepackage{tikz}
\graphicspath{{figures/}{./}}
\usepackage{xcolor}
\usepackage[most]{tcolorbox}
\usepackage{caption}
\usepackage{hyperref}
\hypersetup{colorlinks=true,linkcolor=blue!60!black,citecolor=blue!60!black,urlcolor=blue!60!black}

\theoremstyle{definition}
\newtheorem{definition}{Définition}
\theoremstyle{plain}
\newtheorem{proposition}{Proposition}

% --- encadré "intuition" (vulgarisation) ---
\newtcolorbox{intuition}[1][Intuition]{
  enhanced, breakable,
  colback=blue!4, colframe=blue!55!black, fonttitle=\bfseries,
  title=#1, boxrule=0.6pt, left=4pt,right=4pt,top=3pt,bottom=3pt}

% --- raccourcis mathématiques ---
\newcommand{\R}{\mathbb{R}}
\newcommand{\rhat}{\hat{r}}
\newcommand{\norm}[1]{\left\lVert #1 \right\rVert}
\newcommand{\inner}[2]{\langle #1, #2 \rangle}
\newcommand{\dmodel}{d_{\text{model}}}
\newcommand{\E}{\mathbb{E}}
\DeclareMathOperator{\proj}{proj}
\DeclareMathOperator{\tr}{tr}

\title{\textbf{Model Scanner} — projet \textbf{Jorak} \\[2mm]
\large Détection \emph{reference-free} de l'abliteration des grands modèles de langue\\
\normalsize\emph{(version simplifiée — justification mathématique condensée)}}
\author{Jolan Mocanu}
\date{Version 1.0 --- juillet 2026}

\begin{document}
\maketitle

\begin{abstract}
\noindent
L'\emph{abliteration} retire le comportement de refus d'un grand modèle de langue
(LLM) \textbf{sans réentraînement}, en supprimant une unique direction dans
l'espace interne du modèle. Détecter qu'un modèle a été abliteré est un enjeu de
sûreté et de traçabilité. \textbf{Jorak} est un détecteur \emph{reference-free}
(qui n'a \textbf{pas} besoin du modèle d'origine) reposant sur trois techniques
complémentaires~: la \emph{santé} de l'axe de refus mesurée sur les activations
($d$ de Cohen), une \emph{signature spectrale} des matrices de poids (SVD), et le
\emph{comportement} observable. Ce document condense la justification
mathématique~: on formalise le mécanisme d'abliteration, on définit et on
\emph{commente} chaque métrique, puis on valide le classifieur sur une matrice de
modèles réels dérivés d'une même base (\texttt{Qwen2.5-0.5B}). Le détecteur classe
correctement les trois provenances et n'émet \textbf{aucun faux positif} sur un
modèle simplement fine-tuné. On confirme en outre la signature de poids sur des
modèles abliterés réels diffusés en ligne --- jusqu'à \textbf{8B} et sur des
architectures différentes (Granite-3.2-8B, Qwen3-VL-4B) --- tous deux signalés avec
une confiance maximale.
\end{abstract}

\tableofcontents
\bigskip

% =====================================================================
\section{Introduction}
% =====================================================================

\subsection{Contexte}
Les LLM dits \og alignés\fg{} sont entraînés à \textbf{refuser} certaines requêtes
(dangereuses, illégales, etc.). Arditi et~al.~\cite{arditi2024} ont montré que ce
refus est, dans une large mesure, porté par une \textbf{unique direction} $\rhat$
dans la représentation interne du modèle~: la projection de l'état interne sur
$\rhat$ est élevée pour les requêtes dangereuses, basse pour les inoffensives. On
peut alors \emph{retirer} le refus en supprimant cette direction des poids~: c'est
l'\textbf{abliteration}~\cite{labonne2024}, aujourd'hui automatisée et appliquée à
des centaines de modèles \og uncensored\fg{} diffusés en ligne.

En une phrase~: on ne réentraîne pas le modèle, on \emph{débranche} chirurgicalement
le câble par lequel il exprimait son refus. Le modèle \og sait\fg{} toujours qu'une
requête est dangereuse — sa représentation interne reste intacte — mais il ne peut
plus \emph{agir} sur ce savoir.

\subsection{Problème}
Étant donné un modèle open-source (souvent quantifié, parfois sans \emph{pedigree}),
\textbf{a-t-il été abliteré~?} La difficulté centrale~: en pratique on ne dispose
\emph{pas} du modèle d'origine pour comparer \og avant/après\fg{}. On veut donc un
détecteur \emph{reference-free}. La comparaison avec l'original (\og oracle\fg{}) ne
sert qu'à \emph{valider} la méthode en développement, jamais en production.

\subsection{Jorak~: le projet et ses trois techniques}
\begin{intuition}[Jorak = le projet, pas une métrique]
\textbf{Jorak} est le nom du \emph{projet} (ce scanner et le rapport qu'il produit).
Ce n'est pas une métrique unique, mais un \textbf{ensemble de trois techniques}
de détection complémentaires, qui examinent le modèle sous trois angles~:
\begin{itemize}\itemsep1pt
  \item le \textbf{$d$ de Cohen} — sur les \emph{activations}~: l'axe de refus
  est-il encore \emph{vivant}~? (\S\ref{sec:cohen})
  \item la \textbf{signature spectrale SVD} — sur les \emph{poids}~: le chemin
  d'écriture du refus est-il \emph{sectionné}~? Elle porte trois signaux,
  $\mathcal{A}$, $\mathcal{B}$ et $\mathcal{S}$ (\S\ref{sec:svd}).
  \item le \textbf{comportement} — le modèle \emph{refuse}-t-il encore~?
  (\S\ref{sec:behav})
\end{itemize}
En particulier, le champ \texttt{svd\_alignment} ($\mathcal{A}$) n'est \emph{qu'un}
des signaux de la deuxième technique~: il ne faut pas le confondre avec \og Jorak\fg{}
dans son ensemble.
\end{intuition}

\noindent Ces trois techniques sont combinées par un \textbf{classifieur de
provenance} (\S\ref{sec:classif}) qui range le modèle dans l'une de quatre classes~:
\emph{censuré} / \emph{abliteré} / \emph{fine-tuné-décensuré} / \emph{ambigu}, avec
un score de confiance et les couches touchées.

% =====================================================================
\section{Préliminaires mathématiques}\label{sec:prelim}
% =====================================================================
On note $\dmodel$ la dimension de l'espace interne ($\dmodel=896$ pour
\texttt{Qwen2.5-0.5B}), $L$ le nombre de couches ($L=24$ ici). Les vecteurs sont en
colonne~; $\norm{\cdot}$ est la norme euclidienne, $A^\top$ la transposée, $I$
l'identité.

\subsection{Le \og residual stream\fg{}~: où vivent les représentations}
Un transformeur maintient, à chaque position, un vecteur $h\in\R^{\dmodel}$, le
\textbf{residual stream}~\cite{elhage2021}. Sa propriété clé est qu'il est
\textbf{additif}~: chaque sous-module \emph{lit} $h$, calcule une contribution $c$,
et la \emph{rajoute} ($h\leftarrow h+c$). Deux modules seulement \emph{écrivent}
dans ce flux — la sortie d'attention \texttt{o\_proj} et celle du MLP
\texttt{down\_proj}~; ce sont donc les seules cibles de l'abliteration.

\begin{intuition}[Pourquoi des \og directions\fg{}]
On peut voir $h$ comme un tableau blanc où chaque couche \emph{ajoute} une
annotation sans rien effacer. Comme l'information s'y inscrit de façon largement
\emph{linéaire}, un concept y devient une \textbf{direction} que l'on peut isoler
par projection — l'observation fondatrice de l'interprétabilité
mécaniste~\cite{elhage2021}.
\end{intuition}

\subsection{Projection et cosinus}
Pour $u,v\in\R^{\dmodel}$, le produit scalaire $\inner{u}{v}=u^\top v$ mesure leur
alignement. La \textbf{projection} de $v$ sur une direction unitaire $\rhat$
($\norm{\rhat}=1$) est $\proj_{\rhat}(v)=(\rhat\rhat^\top)v$~; la matrice
$P=\rhat\rhat^\top$ est le \textbf{projecteur} sur la droite de $\rhat$. Son
complément $I-P$ \emph{efface} la composante selon $\rhat$ et garde le reste —
c'est, mot pour mot, l'opération de l'abliteration (\S\ref{sec:meca}). L'angle
entre deux vecteurs se mesure par le \textbf{cosinus}
$\cos(u,v)=\inner{u}{v}/(\norm{u}\,\norm{v})\in[-1,1]$~; le signe d'une direction
étant arbitraire, on travaille souvent avec $\lvert\cos\rvert$.

\subsection{L'hypothèse de représentation linéaire}
On postule que le concept \og la requête est dangereuse\fg{} est encodé
\textbf{linéairement}~: il existe une direction $\rhat$ telle que $\inner{\rhat}{h}$
est élevé pour les exemples \emph{harmful}, bas pour les \emph{harmless}. Cette
idée — des régularités linéaires des plongements de mots~\cite{mikolov2013} à leur
formalisation pour les LLM~\cite{park2023} — transforme une question floue (\og le
modèle est-il censuré~?\fg{}) en une question \emph{géométrique} précise, à laquelle
les outils linéaires suivants répondent un à un.

% =====================================================================
\section{Le mécanisme de l'abliteration}\label{sec:meca}
% =====================================================================
\begin{definition}[Direction de refus]
La \textbf{direction de refus} $\rhat\in\R^{\dmodel}$, $\norm{\rhat}=1$, est la
direction du flux résiduel le long de laquelle le modèle encode et exprime le refus.
\end{definition}

Soit $W$ une matrice qui \emph{écrit} dans le flux résiduel
($W\in\{\texttt{o\_proj},\texttt{down\_proj}\}$). Pour empêcher le modèle d'écrire le
long de $\rhat$, on applique $I-\rhat\rhat^\top$ à chaque sortie, soit
\begin{equation}\label{eq:ablate}
  W' = (I-\rhat\rhat^\top)\,W .
\end{equation}
C'est une perturbation de \textbf{rang $1$}~: le modèle peut tout écrire \emph{sauf}
le long de l'axe du refus.

\begin{proposition}[L'invariant exploité par la détection]\label{prop:invariant}
Après abliteration, $\rhat^\top W'=0$~: la direction de refus appartient au
\textbf{noyau à gauche} de $W'$.
\end{proposition}
\begin{proof}
$\rhat^\top W'=\rhat^\top(I-\rhat\rhat^\top)W=\rhat^\top W-(\rhat^\top\rhat)\rhat^\top W=0$,
car $\rhat^\top\rhat=1$.
\end{proof}

\noindent Cet invariant est \emph{exact} et \emph{permanent} (gravé dans les poids
livrés). L'abliteration ne touche que les poids d'écriture~: la représentation
interne du refus, elle, reste intacte. D'où la \textbf{signature} recherchée~:
\[
  \boxed{\ \text{axe de refus \emph{vivant}}\ \wedge\ \text{chemin d'écriture \emph{sectionné}}\ }
\]
combinaison qu'aucun fine-tuning ne reproduit~: un fine-tuning déplace l'ensemble des
poids et des représentations, il ne \emph{débranche} pas une direction unique.
Détecter revient donc à \emph{retrouver} $\rhat$ — que l'on ne nous donne pas — puis à
\emph{vérifier} qu'elle est à la fois vivante dans les activations et absente des
poids.

% =====================================================================
\section{Les trois techniques de détection de Jorak}\label{sec:outils}
% =====================================================================
Le scanner réalise un \textbf{unique} passage avant (\og scan-once\fg{}) et en dérive
toutes les mesures. On travaille couche par couche~: $h^{h}_{\ell,i}$ (resp.
$h^{hl}_{\ell,i}$) désigne l'état interne au dernier token du prompt \emph{harmful}
(resp. \emph{harmless}) $i$ en sortie de la couche $\ell$, sur $n$ exemples par classe.

\subsection{Prérequis~: la direction de refus par différence de moyennes}
On ne connaît pas $\rhat$~; on l'\emph{estime}. Le séparateur le plus simple de deux
nuages de points est le vecteur qui relie leurs centres de gravité~:
\begin{equation}\label{eq:meandiff}
  \rhat_\ell=\frac{\mu^{h}_\ell-\mu^{hl}_\ell}{\norm{\mu^{h}_\ell-\mu^{hl}_\ell}},
  \qquad
  \mu^{h}_\ell=\tfrac1n\textstyle\sum_i h^{h}_{\ell,i},\quad
  \mu^{hl}_\ell=\tfrac1n\textstyle\sum_i h^{hl}_{\ell,i}.
\end{equation}
C'est la direction \emph{mean-diff}~\cite{arditi2024}~: un séparateur de rang $1$,
sans inversion de matrice, robuste en petite donnée. Elle fournit un \emph{candidat}
d'axe de refus par couche, estimé sans le modèle d'origine. \emph{(Implémenté dans
\texttt{modelscanner/directions/mean\_diff.py}.)}

\subsection{Technique 1 — Santé de l'axe~: le $d$ de Cohen (activations)}\label{sec:cohen}
On projette chaque état sur $\rhat_\ell$ ($p_{\ell,i}=\inner{\rhat_\ell}{h_{\ell,i}}$)
et l'on quantifie la \emph{séparation} des deux nuages de scalaires (harmful vs
harmless) par la \textbf{taille d'effet} de Cohen~\cite{cohen1988}~:
\begin{equation}\label{eq:cohen}
  d_\ell=\frac{\overline{p^{h}_\ell}-\overline{p^{hl}_\ell}}{\sigma_{\text{pool},\ell}},
  \qquad
  \sigma_{\text{pool}}=\sqrt{\tfrac{(n_h-1)s_h^2+(n_{hl}-1)s_{hl}^2}{n_h+n_{hl}-2}} .
\end{equation}

\begin{intuition}[Ce que mesure $d_\ell$]
$d_\ell$ exprime \og de combien d'écarts-types\fg{} les deux classes sont séparées le
long de $\rhat_\ell$. Diviser par $\sigma_{\text{pool}}$ rend la mesure
\emph{sans dimension}, donc comparable d'une couche, d'un modèle et d'une
quantification à l'autre. Conventions usuelles~\cite{cohen1988}~: $0{,}2$ faible,
$0{,}5$ moyen, $0{,}8$ fort. Un $d_\ell$ grand signale un \textbf{axe vivant}~; on
mesure jusqu'à $d_\ell\approx 4{,}5$ dans les couches profondes de
\texttt{Qwen2.5-0.5B}.
\end{intuition}

\noindent\textbf{Garde-fou.} Si l'on calcule $\rhat$ \emph{et} $d$ sur les
\emph{mêmes} sondes, $d$ devient optimiste. Le protocole impose un \textbf{découpage
train/test} des sondes (estimer $\rhat$ sur un sous-ensemble, mesurer $d$ sur un
autre, disjoint). \emph{(Implémenté dans
\texttt{modelscanner/metrics/axis\_health.py}.)} Le $d$ de Cohen mesure la
\emph{vitalité} de l'axe dans les représentations~; il ne dit rien des poids — un
modèle abliteré garde un axe vivant tout en ayant le chemin d'écriture sectionné.
D'où la technique suivante.

\subsection{Technique 2 — Signature spectrale des poids (SVD)}\label{sec:svd}

\paragraph{(a) Lire la direction directement dans la matrice.}
Une première idée serait de projeter \emph{notre} $\rhat$ estimée sur les poids et de
mesurer ce qui \og passe encore\fg{}. Cela \textbf{échoue}~: notre $\rhat_m$
(mean-diff) n'est pas exactement la direction $\rhat_a$ utilisée par l'ablateur,
si bien que $\rhat_m^\top W'\neq 0$ et que le résidu se \emph{noie} dans le bruit
(de l'ordre de $1/\sqrt{\dmodel}\approx 0{,}03$). On cesse donc de \emph{deviner}
$\rhat$ pour la \emph{lire} dans la matrice, grâce à la \textbf{décomposition en
valeurs singulières}~\cite{golub2013}.

\begin{intuition}[La SVD en une image]
$W$ transforme la sphère unité en un \emph{ellipsoïde}. Les longueurs de ses axes
sont les \textbf{valeurs singulières} $\sigma_1\ge\dots\ge\sigma_r\ge 0$, leurs
directions de sortie les \textbf{vecteurs singuliers à gauche} $u_1,\dots,u_r$. Le
plus petit, $u_{\min}$, est la direction que la matrice \og produit le moins\fg{}.
\end{intuition}

\noindent Formellement $W=U\Sigma V^\top$, et $\norm{u_i^\top W}=\sigma_i$.
\textbf{Conséquence (le cœur de la méthode)}~: après abliteration $\rhat_a^\top W'=0$
(Prop.~\ref{prop:invariant}), donc $\rhat_a$ est un vecteur singulier à gauche de
valeur singulière \emph{nulle}~: $u_{\min}(W')\approx\rhat_a$. \emph{On récupère la
direction ablatée sans jamais l'avoir connue.}

\paragraph{(b) Signal $\mathcal{A}$ — alignement inter-couches (le \emph{smoking-gun}).}
Une seule couche ne suffit pas (toute matrice a un $u_{\min}$). La signature décisive
est que, sur un modèle abliteré, \emph{toutes} les couches partagent la \emph{même}
$u_{\min}\approx\rhat_a$ (la même direction ayant été retirée partout). On empile les
$u_{\min}$ en $\mathbf{U}\in\R^{L\times\dmodel}$ (lignes unitaires) et l'on mesure leur
cohérence~:
\begin{equation}\label{eq:align}
  \mathcal{A}=\frac{\sigma_1(\mathbf{U})}{\norm{\mathbf{U}}_F}\in\Big[\tfrac{1}{\sqrt{L}},\,1\Big].
\end{equation}

\begin{intuition}[Lire $\mathcal{A}$]
Si toutes les couches partagent la même direction, $\mathbf{U}$ est de rang $1$ et
$\mathcal{A}\approx 1$. Si les directions sont sans rapport (bruit non corrélé d'un
modèle sain), $\mathcal{A}\approx 1/\sqrt{L}$. Avec $L=24$, cela donne $\approx 0{,}20$
(sain) contre $\approx 0{,}90$ (abliteré)~: une séparation franche, \emph{indépendante
de toute sonde}, calculable sur CPU. On garde aussi un \textbf{signal par couche}
$a_\ell=\lvert\cos(u_{\min}^{(\ell)},w)\rvert$ ($w$~=~axe partagé dominant) pour savoir
\emph{quelles} couches ont été touchées (lignes de la heatmap).
\end{intuition}

\paragraph{(c) Signal $\mathcal{B}$ — alignement de \emph{bande} (ablation localisée).}
$\mathcal{A}$ moyenne sur \emph{toutes} les couches. Or une ablation \emph{furtive}
(type Heretic, qui optimise pour préserver les capacités) ne touche qu'une bande de
couches~: la moyenne globale est alors \emph{diluée} et repasse sous le seuil. On
ajoute donc une statistique \emph{par bande}~:
\begin{equation}\label{eq:band}
  \mathcal{B}=\max_{0\le i\le L-w}\ \frac1w\sum_{\ell=i}^{i+w-1} a_\ell ,
  \qquad w\approx L/6 ,
\end{equation}
la moyenne maximale de l'alignement par couche sur une fenêtre de $w$ couches
\emph{consécutives}. $\mathcal{B}$ est \emph{agnostique en position}~: un modèle sain
(bruit dispersé) ne forme aucune bande ($\mathcal{B}$ bas)~; une ablation localisée
(bande contiguë) donne $\mathcal{B}$ élevé, où qu'elle se trouve.

\paragraph{(d) Signal $\mathcal{S}$ — alignement de \emph{sous-espace} (multi-direction).}
$\mathcal{A}$ et $\mathcal{B}$ supposent une \emph{unique} direction retirée. Or,
empiriquement, le refus vit dans \emph{plusieurs} directions indépendantes, voire un
\emph{cône} de concepts~\cite{wollschlager2025,piras2025}, et les abliterateurs
récents (Gabliteration~\cite{gabliteration}) retirent un
\emph{sous-espace} de $k>1$ directions~: la trace est alors \emph{étalée} sur les $k$
plus petites valeurs singulières et aucun $u_{\min}$ unique n'est partagé. On
\emph{généralise} donc au sous-espace de bas rang. Soit $B_\ell\in\R^{\dmodel\times k}$
les $k$ vecteurs singuliers à gauche de plus petites valeurs singulières de la couche
$\ell$~; on cherche le sous-espace $k$-D le plus \emph{partagé}~:
\begin{equation}\label{eq:subspace}
  M=\sum_{\ell=1}^{L} B_\ell B_\ell^\top,\qquad
  \mathcal{S}=\frac{1}{Lk}\sum_{i=1}^{k}\lambda_i(M)\ \in\Big[\tfrac{k}{\dmodel},\,1\Big],
\end{equation}
où $\lambda_1\ge\dots$ sont les valeurs propres de $M$. $\mathcal{S}\to 1$ si toutes
les couches partagent le même sous-espace $k$-D (abliteré multi-direction)~;
$\mathcal{S}\approx k/\dmodel$ si les sous-espaces sont sans rapport (sain). Pour
$k=1$, $\mathcal{S}=\mathcal{A}^2$~: c'est bien une généralisation.

\begin{intuition}[Les trois signaux de poids, en un coup d'œil]
$\mathcal{A}$ détecte l'ablation \emph{totale} (toutes les couches)~; $\mathcal{B}$
l'ablation \emph{localisée} (une bande de couches)~; $\mathcal{S}$ l'ablation
\emph{multi-direction} (un sous-espace de $k$ directions). Le critère de
\og poids sectionnés\fg{} est~: \textbf{$\mathcal{A}$, $\mathcal{B}$ \emph{ou}
$\mathcal{S}\ge 0{,}5$}. \emph{(Tous calculés en NumPy pur, sans inférence, dans
\texttt{modelscanner/metrics/jorak.py} — le cœur \emph{reference-free} et CPU de Jorak.)}
\end{intuition}

\subsection{Technique 3 — Le comportement (behavioral)}\label{sec:behav}
On fait \emph{générer} une réponse à des prompts harmful et l'on mesure le
\textbf{taux de refus} par appariement de marqueurs (\og I cannot\fg{}, \og désolé\fg{},
\dots) — l'heuristique standard d'évaluation~\cite{zou2023}. C'est le \emph{seul}
signal qui distingue un modèle qui \emph{refuse} d'un modèle qui \emph{obéit}.

\begin{intuition}[Pourquoi le comportement est indispensable]
Un modèle simplement \emph{fine-tuné-décensuré} garde, lui aussi, un axe \emph{vivant}
($d$ élevé) \emph{et} des poids \emph{intacts} ($\mathcal{A},\mathcal{B},\mathcal{S}$
bas)~: pour les seules activations et poids, il est \textbf{indistinguable} d'un
censuré. Seul le \emph{comportement} les sépare — le censuré refuse, le fine-tuné
obéit. Le plan comportemental n'est donc pas optionnel.
\end{intuition}

\noindent On mesure de plus le \textbf{sur-refus du harmless} (refus de requêtes
\emph{inoffensives}) et l'\textbf{asymétrie multilingue} (écart de refus entre
langues)~: deux marqueurs de \emph{dégât}. Un modèle sainement fine-tuné y reste
cohérent~; une ablation, elle, abîme le comportement (sur-refus, refus inégal selon la
langue). \emph{(Implémenté dans \texttt{modelscanner/metrics/behavioral.py}.)}

\subsection{Le classifieur de provenance}\label{sec:classif}
On combine les indicateurs scalaires (seuils calibrés)~: \emph{axe vivant}
($\max_\ell d_\ell\ge 1$), \emph{poids sectionnés} ($\mathcal{A}$, $\mathcal{B}$
\textbf{ou} $\mathcal{S}\ge 0{,}5$), \emph{obéit} (taux de refus $<0{,}5$). La règle~:
\[
\begin{array}{l@{\quad}l}
  \text{poids sectionnés} & \Rightarrow \textbf{ABLITÉRÉ}\ (\text{axe vivant attendu})\\
  \text{sinon, refuse} & \Rightarrow \textbf{CENSURÉ}\\
  \text{sinon, obéit} & \Rightarrow \textbf{FINE-TUNÉ-DÉCENSURÉ}\\
  \text{incohérent} & \Rightarrow \textbf{AMBIGU}
\end{array}
\]
On visualise la règle dans un \textbf{plan de provenance}~: en abscisse l'alignement
$\mathcal{A}$ (poids sectionnés $\rightarrow$), en ordonnée le taux de refus
($\uparrow$ refuse). La sévérance des poids ($x$) isole l'abliteré~; le comportement
($y$) sépare censuré et fine-tuné. La \textbf{confiance} rendue est la \emph{marge}
du signal décisif (sa distance au seuil).

\paragraph{Sur les seuils.} Les seuils de décision sont des \emph{valeurs par défaut
raisonnées}, non ajustées au banc d'essai. Pour les signaux de poids, $0{,}5$ est le
milieu entre le plancher théorique d'un modèle décorrélé ($1/\sqrt{L}\approx 0{,}2$
pour $L=24$) et le plafond de rang $1$ ($1$)~; pour la santé d'axe, $d\ge 1$ suit la
convention \og effet fort\fg{} de Cohen~\cite{cohen1988}~; pour le comportement,
refus~$<0{,}5$ est le partage majoritaire naturel. Un balayage systématique par signal
(ROC) est laissé aux travaux futurs --- les marges rapportées en
\S\ref{sec:resultats} montrent que les défauts actuels sont loin de la frontière de
décision.

\paragraph{Garde-fou anti-faux-positif (veto comportemental).} La règle ci-dessus
admet une exception cruciale, issue de la campagne $8$B~: si les poids paraissent
sectionnés \emph{mais} que le modèle \textbf{refuse encore largement} (taux de refus
$\ge 0{,}5$), le \emph{comportement prime} et le verdict reste \textbf{CENSURÉ}. Un
signal de poids peut en effet être un \emph{faux positif} (p.~ex.\ une bande
naturellement alignée sur un modèle de base)~; un abliteré réussi, lui, \emph{obéit}
par construction et ne déclenche jamais ce veto.

\begin{intuition}[Le cas furtif, en bref]
Une ablation \emph{multi-direction norm-preserving} peut laisser les poids
\og intacts\fg{} aux yeux de la SVD ($\mathcal{A},\mathcal{B},\mathcal{S}$ sous le
seuil). Le \emph{comportement} la trahit alors~: un axe \emph{fortement} vivant
combiné à l'obéissance et à un dégât incohérent (sur-refus / asymétrie multilingue)
ne s'explique pas par un fine-tuning propre. La confiance rendue est plus
\emph{modérée}, car la preuve est \emph{indirecte}.
\end{intuition}

\noindent\emph{(Implémenté dans \texttt{modelscanner/classifier/scanner.py}~:
\texttt{classify(...)} est une fonction pure, testable hors modèle~; \texttt{scan(model)}
orchestre le scan-once et assemble un \texttt{ScanResult}.)}

% =====================================================================
\section{Résultats}\label{sec:resultats}
% =====================================================================
\textbf{Banc d'essai.} Trois modèles dérivés de la \emph{même} base
\texttt{Qwen2.5-0.5B} (même architecture, même taille — comparaison propre)~:
\texttt{Qwen2.5-0.5B-Instruct} (\emph{censuré}),
\texttt{huihui-ai/\dots-abliterated} (\emph{abliteré} réel),
\texttt{dphn/Dolphin3.0-Qwen2.5-0.5B} (\emph{fine-tuné-décensuré}, le faux positif à
éviter), plus un Heretic réel (ablation \emph{localisée}). \textbf{Sondes}~: $40$ paires
harmful/harmless appariées en longueur et registre~; split train/test des sondes.

\begin{table}[H]
\centering
\caption{Verdict du scanner sur la matrice. $\mathcal{A}$~: alignement SVD global~;
$\mathcal{B}$~: alignement de bande (localisé)~; $\mathcal{S}$~: alignement de
sous-espace ($k{=}4$)~; $d_{\max}$~: santé d'axe~; \og refus\fg{}~: taux de refus.}
\label{tab:matrice}
\begin{tabular}{llccccc}
\toprule
Modèle & Label & $d_{\max}$ & $\mathcal{A}$ & $\mathcal{B}$ & $\mathcal{S}$ & refus \\
\midrule
\texttt{Qwen2.5-0.5B-Instruct} & \textsc{censuré} & $4{,}5$ & $0{,}24$ & $0{,}28$ & $0{,}08$ & $0{,}88$ \\
\texttt{huihui\dots abliterated} & \textsc{abliteré} & $2{,}6$ & $\mathbf{0{,}90}$ & $\mathbf{0{,}99}$ & $0{,}29$ & $0{,}06$ \\
\texttt{grisun0\dots heretic} & \textsc{abliteré} & $3{,}3$ & $0{,}29$ & $\mathbf{0{,}67}$ & $0{,}10$ & --- \\
\texttt{Dolphin3.0-Qwen2.5-0.5B} & \textsc{fine-tuné} & $3{,}0$ & $0{,}24$ & $0{,}28$ & $0{,}08$ & $0{,}06$ \\
\bottomrule
\end{tabular}
\end{table}

\noindent\emph{\small Reproduire~: \texttt{modelscanner scan <modèle> --behavioral
--report}~; chaque valeur est lue telle quelle dans le \texttt{scan\_log.json} émis
(jeu de sondes~: $40$ paires harmful/harmless, split train/test).}

\noindent Les quatre provenances sont \textbf{correctement classées}, et le modèle
fine-tuné n'est \textbf{pas} pris pour un abliteré (faux positif critique évité).
Trois lectures résument la matrice~:
\begin{itemize}\itemsep2pt
  \item \textbf{La SVD est le vrai smoking-gun.} Seul $\mathcal{A}$ sépare l'abliteré
  ($0{,}90$) du reste ($0{,}24$). La projection naïve de $\rhat$ sur les poids vaut
  $\approx 0{,}038$ \emph{pour les trois} modèles (= le plancher aléatoire $0{,}033$)~:
  elle ne voit rien. C'est la signature \emph{spectrale} qui tranche.
  \item \textbf{Le comportement sépare ce que les poids ne séparent pas.} Le fine-tuné
  a un axe vivant ($d_{\max}=3{,}0$) et des poids intacts ($\mathcal{A}=0{,}24$)~:
  identique au censuré en reference-free. Seul le refus les distingue ($0{,}06$ vs
  $0{,}88$).
  \item \textbf{La bande rattrape l'ablation localisée.} Le modèle Heretic a
  $\mathcal{A}=0{,}29$ (sous le seuil~: l'ablation serait \emph{ratée} par le global),
  mais $\mathcal{B}=0{,}67$~: Heretic n'a touché qu'une bande \emph{profonde}
  (couches $19$--$22$ sur $24$), révélée par la fenêtre glissante.
\end{itemize}
La colonne $\mathcal{S}$ reste \emph{dormante} ($\le 0{,}29$) sur ces ablations
\emph{mono}-direction~: par construction, elle ne se déclenche que sur le
\emph{multi}-direction, et n'introduit ici aucun faux positif.

\begin{figure}[H]
\centering
\includegraphics[width=0.72\linewidth]{scatter.png}
\caption{\textbf{Plan de provenance.} Abscisse~: alignement SVD $\mathcal{A}$ (poids
sectionnés $\rightarrow$). Ordonnée~: taux de refus ($\uparrow$ refuse). Les trois
modèles occupent trois quadrants distincts, correspondant exactement à leur provenance.}
\label{fig:scatter}
\end{figure}

\begin{figure}[H]
\centering
\includegraphics[width=0.9\linewidth]{heatmap-ablated.png}
\caption{\textbf{Signature par couche du modèle abliteré.} La ligne
\emph{alignement poids par couche} $a_\ell$ (échelle $[0,1]$) est \textbf{saturée} sur
presque toutes les couches~: l'ablation a touché tout le réseau.}
\label{fig:heatmap-abl}
\end{figure}

\begin{figure}[H]
\centering
\includegraphics[width=0.9\linewidth]{heatmap-censored.png}
\caption{\textbf{Modèle censuré (témoin).} Sur la même échelle $[0,1]$, la ligne
\emph{alignement poids par couche} reste \textbf{sombre} ($\approx 0{,}2$)~: aucune
direction partagée n'a été retirée.}
\label{fig:heatmap-cen}
\end{figure}

\noindent\textbf{Confiance.} La confiance du verdict est la \emph{marge} du signal
décisif~: $0{,}81$ (abliteré), $0{,}75$ (censuré), $0{,}88$ (fine-tuné).

\subsection{Validation terrain à l'échelle et sur d'autres architectures}
Au-delà du banc contrôlé $0{,}5$B, on a passé Jorak sur des modèles abliterés réels
diffusés en ligne, à plus grande échelle et sur des architectures différentes
(Table~\ref{tab:field}). Les deux sont correctement signalés \textsc{abliteré} avec une
\textbf{confiance maximale}~; pour Granite, le verdict correspond à la provenance
déclarée (\emph{attendu}~$\checkmark$).

\begin{table}[H]\centering
\caption{Validation terrain sur des modèles abliterés réels (lignée huihui). Valeurs
lues dans le \texttt{scan\_log.json} de chaque modèle.}
\label{tab:field}
\begin{tabular}{llccccc}
\toprule
Modèle & Arch\,·\,taille & $\mathcal{A}$ & $\mathcal{B}$ & $\mathcal{S}_{k=4}$ & $d_{\max}$ & refus \\
\midrule
\texttt{Qwen3-VL-4B-abliterated} & qwen3\_vl · 4B & $1{,}00$ & $1{,}00$ & $0{,}39$ & $4{,}97$ & $0{,}00$ \\
\texttt{granite-3.2-8B-abliterated} & granite · 8B & $0{,}99$ & $1{,}00$ & $0{,}50$ & $5{,}07$ & $0{,}06$ \\
\bottomrule
\end{tabular}
\end{table}

\noindent Trois enseignements~: \textbf{(i)} la signature inter-couches $\mathcal{A}$
\emph{sature} ($0{,}99$--$1{,}00$) à $8$B et sur un modèle vision--langage --- ce n'est
pas un artefact du $0{,}5$B~; \textbf{(ii)} la combinaison-signature --- axe
\emph{vivant} ($d_{\max}\approx 5$), chemin d'écriture \emph{sectionné} et
\emph{obéissance} (refus~$\approx 0$) --- tient à l'échelle~; \textbf{(iii)} les
capacités restent largement intactes (Granite~: exactitude factuelle $0{,}88$,
perplexité $11{,}2$), soit une ablation \emph{propre} que le plan des poids détecte
malgré tout d'emblée.

\begin{figure}[H]\centering
\begin{tikzpicture}[yscale=0.6]
\foreach \name/\y/\a/\b in {%
  Heretic (0.5B)/2/0.79/0.92,
  Qwen3-VL (4B)/1/0.28/0.42,
  Granite (8B)/0/0.18/0.33}{
  \draw[gray!45] (0,\y) rectangle (8,\y+0.6);
  \fill[blue!55!black] (8*\a,\y) rectangle (8*\b,\y+0.6);
  \node[left] at (-0.15,\y+0.3) {\small \name};
}
\node[below] at (0,-0.05) {\scriptsize entrée};
\node[below] at (8,-0.05) {\scriptsize sortie};
\node[below=7pt] at (4,-0.05) {\scriptsize profondeur de couche (normalisée) $\rightarrow$};
\end{tikzpicture}
\caption{\textbf{Le signal de bande $\mathcal{B}$ est agnostique en position.} La bande
contiguë de couches sectionnées (grisée) se situe à une profondeur \emph{différente}
selon le modèle --- profonde pour Heretic, médiane pour Qwen3-VL et Granite --- mais la
fenêtre glissante (éq.~\eqref{eq:band}) la rattrape dans tous les cas.}
\label{fig:band}
\end{figure}

\paragraph{Portée honnête.} Ces deux modèles sont des \emph{positifs} (tous deux
ablatés). Une matrice $8$B complète avec \emph{témoins négatifs} (un censuré et un
fine-tuné-seul à $8$B) est laissée aux travaux futurs~; la propriété
\emph{aucun-faux-positif} est donc établie sur la matrice $0{,}5$B
(\S\ref{sec:resultats}) et pas encore à $8$B. L'échantillon ne couvre en outre qu'une
seule lignée/outil (huihui).

\paragraph{Leçon méthodologique (anti-faux-positif).} La cible \texttt{down\_proj}
(sortie du MLP) a une bande $\mathcal{B}$ \emph{naturellement haute sur les modèles de
base} (non discriminante) et a produit des faux positifs lors d'un premier essai
multi-cible~; le verdict de campagne s'appuie donc sur \texttt{o\_proj} seul, complété
par le \textbf{veto comportemental} (\S\ref{sec:classif}). Le journal de scan trace
désormais la \emph{ventilation par cible} ($\mathcal{A}$, $\mathcal{B}$ par
\texttt{o\_proj}/\texttt{down\_proj} séparément) pour rendre tout réexamen non ambigu.

% =====================================================================
\section{Travaux connexes}\label{sec:related}
% =====================================================================
Jorak se situe au croisement de trois courants récents~: la \emph{géométrie} du refus,
les \emph{défenses} contre l'abliteration, et le \emph{paysage} des outils d'ablation.

\paragraph{Du mécanisme unique aux directions multiples.}
Le point de départ — une \emph{unique} direction de refus~\cite{arditi2024} — a été
nuancé. Wollschläger et~al.~\cite{wollschlager2025} montrent, par une approche à base
de gradient, l'existence de \emph{plusieurs} directions mécaniquement indépendantes,
voire de \emph{cônes de concepts} multi-dimensionnels, et que l'orthogonalité
n'implique pas l'indépendance sous intervention. Piras et~al.~\cite{piras2025}
extraient de même plusieurs directions par cartes auto-organisatrices (SOM) et
établissent qu'ablater \emph{plusieurs} directions supprime le refus plus efficacement
qu'une seule (leur méthode \emph{généralise} d'ailleurs la différence de moyennes de
l'éq.~\eqref{eq:meandiff}). Zhang \& Sun~\cite{zhang2025dbdi} scindent enfin la
direction unique en deux processus — \emph{détection du danger} et \emph{exécution du
refus}. Ces résultats justifient directement notre signal de \emph{sous-espace}
$\mathcal{S}$ (\S\ref{sec:svd}d)~: on ne peut supposer qu'un seul $u_{\min}$ porte toute
la trace de l'ablation.

\paragraph{Attaque, défense, détection.}
L'abliteration a suscité ses \emph{défenses}. Abu Shairah et~al.~\cite{abushairah2025}
proposent un fine-tuning \emph{extended-refusal} (le modèle \emph{justifie} son refus)
qui limite la chute du taux de refus à $10\%$ au plus après attaque, contre
$70$--$80\%$ pour la baseline. Jorak est \emph{complémentaire}~: ce n'est pas une
défense mais un outil de \emph{détection}/forensique. À noter qu'un modèle ainsi
défendu \emph{refuse encore} après tentative d'ablation — précisément le cas que notre
\emph{veto comportemental} (\S\ref{sec:classif}) est conçu pour ne pas confondre avec
un modèle sain.

\paragraph{Le paysage des outils d'ablation.}
Young~\cite{young2025} compare quatre outils d'abliteration (Heretic, DECCP, ErisForge,
FailSpy) sur seize modèles instruct ($7$--$14$B) et confirme que le dégât de capacités
dépend de l'outil, le raisonnement mathématique (GSM8K) étant le plus sensible — ce qui
conforte à la fois notre campagne multi-outils (\S\ref{sec:resultats}) et nos
\emph{marqueurs de dégât} (sur-refus, asymétrie). Enfin, les méthodes
\emph{norm-preserving} récentes~\cite{dang2026,grimjim}, qui remplacent l'ablation
\og binaire\fg{} par une rotation préservant la norme, constituent exactement la classe
d'ablation \emph{furtive} que nous listons comme limite ouverte (\S\ref{sec:robustesse}).

% =====================================================================
\section{Robustesse et limites}\label{sec:robustesse}
% =====================================================================
\begin{itemize}\itemsep2pt
  \item \textbf{Quantification.} $\mathcal{A}$ est une signature \emph{inter-couches}~;
  le bruit de quantification étant largement décorrélé d'une couche à l'autre, il
  dégrade peu une cohérence \emph{partagée}. On conjecture la survie de la signature
  jusqu'à $4$ bits (seuil glissant, \og quant-aware\fg{}) — à valider sur \texttt{bf16},
  \texttt{bnb-4bit}, GGUF.
  \item \textbf{Sophistication de l'ablation.} L'ablation localisée est couverte par
  $\mathcal{B}$~; le multi-direction par $\mathcal{S}$ et, à défaut, par le relais
  comportemental. Restent ouverts~: la projection \emph{partielle} ($\alpha<1$, qui fait
  décroître $\mathcal{A},\mathcal{B},\mathcal{S}$ continûment) et le norm-preserving
  \emph{par lignes}~\cite{dang2026,grimjim}, qui \emph{déplace} le noyau gauche et atténue
  $\mathcal{S}$.
  \item \textbf{Architectures et formats.} Les modèles MoE / MLA demandent un
  traitement par expert~; pour les formats GGUF non exécutables, le plan d'activations
  est indisponible (repli sur poids $+$ comportemental).
\end{itemize}

% =====================================================================
\section{Conclusion}
% =====================================================================
Le fil rouge est simple~: \emph{le refus est une direction}~; l'\emph{abliteration
efface cette direction des poids d'écriture} en laissant un invariant exact
($\rhat^\top W'=0$)~; cet invariant se \emph{lit} dans le spectre des matrices, sans
connaître la direction ni le modèle d'origine, via l'\textbf{alignement des plus
petits vecteurs singuliers}. Jorak combine trois techniques — la \emph{santé de l'axe}
($d$ de Cohen), la \emph{signature spectrale} des poids ($\mathcal{A},\mathcal{B},
\mathcal{S}$) et le \emph{comportement} — pour situer un modèle parmi quatre
provenances. Sur une matrice de modèles réels, le détecteur classe correctement les
trois provenances et évite le faux positif critique. Toute la démarche est
implémentée et reproductible dans la librairie \texttt{modelscanner}.

% =====================================================================
\begin{thebibliography}{99}
\bibitem{arditi2024} A.~Arditi \emph{et al.}, \emph{Refusal in Language Models Is
Mediated by a Single Direction}, arXiv:2406.11717, 2024.
\bibitem{labonne2024} M.~Labonne, \emph{Uncensor any LLM with abliteration},
Hugging Face Blog, 2024.
\bibitem{elhage2021} N.~Elhage \emph{et al.}, \emph{A Mathematical Framework for
Transformer Circuits}, Transformer Circuits Thread (Anthropic), 2021.
\bibitem{mikolov2013} T.~Mikolov, W.~Yih, G.~Zweig, \emph{Linguistic Regularities
in Continuous Space Word Representations}, NAACL-HLT, 2013.
\bibitem{park2023} K.~Park, Y.~J.~Choe, V.~Veitch, \emph{The Linear Representation
Hypothesis and the Geometry of Large Language Models}, arXiv:2311.03658, 2023.
\bibitem{cohen1988} J.~Cohen, \emph{Statistical Power Analysis for the Behavioral
Sciences}, 2\textsuperscript{e} éd., Lawrence Erlbaum, 1988.
\bibitem{golub2013} G.~H.~Golub, C.~F.~Van~Loan, \emph{Matrix Computations},
4\textsuperscript{e} éd., Johns Hopkins University Press, 2013.
\bibitem{zou2023} A.~Zou \emph{et al.}, \emph{Universal and Transferable Adversarial
Attacks on Aligned Language Models}, arXiv:2307.15043, 2023.
\bibitem{gabliteration} G.~Gülmez, \emph{Gabliteration: Adaptive Multi-Directional
Neural Weight Modification for Selective Behavioral Alteration in Large Language
Models}, arXiv:2512.18901, 2025.
\bibitem{grimjim} grimjim, \emph{Norm-Preserving Biprojected Abliteration},
Hugging Face, 2025.
\bibitem{wollschlager2025} T.~Wollschläger, J.~Elstner, S.~Geisler, V.~Cohen-Addad,
S.~Günnemann, J.~Gasteiger, \emph{The Geometry of Refusal in Large Language Models:
Concept Cones and Representational Independence}, arXiv:2502.17420, 2025.
\bibitem{piras2025} G.~Piras, R.~Mura, F.~Brau, L.~Oneto, F.~Roli, B.~Biggio,
\emph{SOM Directions are Better than One: Multi-Directional Refusal Suppression in
Language Models}, arXiv:2511.08379, 2025.
\bibitem{zhang2025dbdi} P.~Zhang, P.~Sun, \emph{Differentiated Directional
Intervention: A Framework for Evading LLM Safety Alignment}, arXiv:2511.06852, 2025.
\bibitem{abushairah2025} H.~Abu~Shairah, H.~A.~K.~Hammoud, B.~Ghanem, G.~Turkiyyah,
\emph{An Embarrassingly Simple Defense Against LLM Abliteration Attacks},
arXiv:2505.19056, 2025.
\bibitem{young2025} R.~J.~Young, \emph{Comparative Analysis of LLM Abliteration
Methods: A Cross-Architecture Evaluation}, arXiv:2512.13655, 2025.
\bibitem{dang2026} Q.-A.~Dang, C.~Ngo, \emph{Selective Steering: Norm-Preserving
Control Through Discriminative Layer Selection}, arXiv:2601.19375, 2026.
\end{thebibliography}

\end{document}
